\section{\texorpdfstring{\texttt{Prop}: the type of statements}{Prop: the type of statements}}\label{sec:04-propositions-and-proofs:01-prop:1}

Chapter 3 closed with structures bundling data and proofs side by side, without yet saying what a ``proof,'' as a piece of Lean data, actually is. That gap closes here. Chapters 1 and 3 built terms and types for ordinary mathematical objects, numbers, points, groups-to-be, and this chapter turns the same machinery on \emph{statements} and their proofs. Alongside \texttt{Type}, Lean has \texttt{Prop}, the type of logical propositions. A term of type \texttt{P : Prop} is a \textbf{proof} of \texttt{P}. This is the \textbf{Curry--Howard correspondence}. Propositions are types, and proofs are programs.

\subsection{The Curry--Howard correspondence, in full}\label{sec:04-propositions-and-proofs:01-prop:2}

That one-line slogan is easy to state, but its importance is easy to miss on first read. What follows is the full dictionary it stands for, a two-way correspondence between logical connectives and type formers. This book uses each row repeatedly, starting in the next few sections.

\begin{longtable}[]{@{}
  >{\raggedright\arraybackslash}p{(\columnwidth - 4\tabcolsep) * \real{0.3333}}
  >{\raggedright\arraybackslash}p{(\columnwidth - 4\tabcolsep) * \real{0.3333}}
  >{\raggedright\arraybackslash}p{(\columnwidth - 4\tabcolsep) * \real{0.3333}}@{}}
\toprule\noalign{}
\begin{minipage}[b]{\linewidth}\raggedright
Logic
\end{minipage} & \begin{minipage}[b]{\linewidth}\raggedright
Type theory
\end{minipage} & \begin{minipage}[b]{\linewidth}\raggedright
Lean notation
\end{minipage} \\
\midrule\noalign{}
\endhead
\bottomrule\noalign{}
\endlastfoot
proposition \(P\) & type \(P\) & \texttt{P : Prop} \\
proof of \(P\) & term of type \(P\) & \texttt{p : P} \\
\(P\) implies \(Q\) & function type & \texttt{P → Q} \\
\(P\) and \(Q\) & product type & \texttt{P ∧ Q} \\
\(P\) or \(Q\) & sum (coproduct) type & \texttt{P ∨ Q} \\
false & empty type (no constructors) & \texttt{False} \\
not \(P\) & function type to the empty type & \texttt{¬P} (:= \texttt{P → False}) \\
for all \(x\), \(P(x)\) & dependent function (Π-) type & \texttt{∀ x, P x} \\
there exists \(x\) with \(P(x)\) & dependent pair (Σ-)type & \texttt{∃ x, P x} \\
proof by cases on a disjunction & pattern match / \href{https://lean-lang.org/doc/reference/latest/Tactic-Proofs/Tactic-Reference/}{\texttt{cases}} & \texttt{Or.elim}, \texttt{cases h with ...} \\
a direct proof (construction) & a term built from constructors & \texttt{⟨\_\allowbreak{}, \_\allowbreak{}⟩}, \texttt{Or.inl \_\allowbreak{}}, \texttt{fun x => \_\allowbreak{}} \\
\end{longtable}

Consider a few rows concretely. ``\(P\) and \(Q\)'' corresponds to a \emph{product} type because a proof of \(P \wedge Q\) is genuinely a \emph{pair}, a proof of \(P\) together with a proof of \(Q\). This is exactly the \texttt{⟨hp, hq⟩} seen in the next section. ``\(P\) or \(Q\)'' corresponds to a \emph{sum} type because a proof of \(P \vee Q\) is a \emph{choice}, either a proof of \(P\) (tagged \texttt{Or.inl}) or a proof of \(Q\) (tagged \texttt{Or.inr}), never both, and never neither. ``Not \(P\)'' being \texttt{P → False} says that a proof that \(P\) is false is a \emph{procedure} that would turn any (hypothetical) proof of \(P\) into a proof of the impossible proposition \texttt{False}. In other words, it is a witness that no such proof of \(P\) could exist.

The correspondence goes deeper than just matching up connectives with type formers, though. It also extends to \emph{proofs themselves} behaving like \emph{programs}. Simplifying a proof (removing a detour, such as proving \(P \wedge Q\) and then immediately taking the left projection to recover a proof of \(P\)) corresponds exactly to a program taking a computation step (here, β-reduction eliminating a constructor immediately followed by the matching projection). This is why the tactics of Chapter 5, which \emph{build} proof terms, and the discussion of reduction and definitional equality in Chapter 6, are really talking about one and the same underlying process, just seen from two angles. Proof simplification and program evaluation are the same operation, described differently depending on whether the term is regarded as ``a proof'' or ``a computation.''

\begin{quote}
Read more. If ``propositional logic'' or ``natural deduction'' above are not already familiar, \hyperref[sec:04-propositions-and-proofs:02-logic-recap:1]{the next section} recaps standard logic from scratch, with no Lean involved, before this correspondence gets applied to it. \hyperref[sec:02-terminology-and-coc:02-pi-sigma-and-coc:1]{Chapter 2, Section 2} makes the correspondence fully precise, extending it down to the untyped λ-calculus underneath both proofs and ordinary functions. The progress and preservation theorems of \hyperref[sec:06-rigor-check:03-typing-rules-and-safety:1]{Chapter 6, Section 3} are the formal statement of ``a proof never reduces to something of the wrong type,'' i.e.~``well-typed proofs do not go wrong.''
\end{quote}

\begin{lstlisting}
#check (2 + 2 = 4)     -- 2 + 2 = 4 : Prop

example : 2 + 2 = 4 := rfl
\end{lstlisting}

\href{https://lean-lang.org/doc/reference/latest/Tactic-Proofs/Tactic-Reference/}{\texttt{rfl}} is the proof ``both sides compute to the same thing'' (\textbf{refl}exivity). \texttt{example} states a proposition and immediately supplies a proof (an anonymous, unnamed \texttt{theorem}).

\begin{mathreading}

The Curry--Howard correspondence identifies a proposition \(P\) with the \emph{set of its proofs}. \(P\) is true exactly when that set is nonempty, i.e.~when there exists some term \(p : P\) (``\(P\) is inhabited''). This is why \texttt{Prop} behaves like a truth value rather than an ordinary type. Every proof-set of a proposition is either empty (false) or, up to proof irrelevance, has exactly one element (true), with no room for a ``third'' proof genuinely different from the rest. Proving \(P\) is exactly exhibiting an element \(p \in P\); nothing more is meant by ``provable.'' The proof \texttt{rfl : 2 + 2 = 4} is the reflexivity witness \(\mathrm{refl}_4\) of the equality relation, valid precisely because both sides reduce to the same normal form \(4\). This is equality of terms that are \emph{definitionally} equal, the strictest notion of ``\(=\)''.

\end{mathreading}

\subsection{Sources, quoted}\label{sec:04-propositions-and-proofs:01-prop:3}

Formal definitions and citations for this section, gathered here for reference (full entries in the \hyperref[sec:root:bibliography:1]{Bibliography}):

\begin{itemize}
\tightlist
\item
  \textbf{Curry--Howard correspondence.} ``The fact that the rules for implication in a proof system for natural deduction correspond exactly to the rules governing abstraction and application for functions is an instance of the Curry-Howard isomorphism, sometimes known as the propositions-as-types paradigm'' (\cite{TPIL4}, ``Propositions and Proofs''; the correspondence traces to a 1969 manuscript by Howard, circulated privately and formally published as \cite{Howard1980}). The name carries two people because the observation is older than the paper by Howard. Curry noticed the correspondence between implicational formulas and combinator types in 1934, and Howard extended and sharpened it in 1969 (Wadler, ``Propositions as Types,'' \emph{CACM} 58(12), 2015: ``a correspondence observed by Curry in 1934 and refined by Howard in 1969''). Picture it like this. A mathematical claim is a blueprint, and a valid proof is an actual working machine built to that blueprint. There is no separate ``proof of correctness'' apart from having successfully built the machine, which is exactly why, in Lean, propositions are types and proofs are the terms (the ``machines'') that inhabit them.
\item
  Howard (\cite{Howard1980}) is the original source of the correspondence this section is named for. Per Sørensen \& Urzyczyn, \emph{Lectures on the Curry-Howard Isomorphism}, Studies in Logic and the Foundations of Mathematics vol.~149, Elsevier, 2006 (a secondary source corroborating this history, not the paper by Howard itself), the manuscript by Howard was ``privately circulated'' from 1969 and not formally published until 1980, in a Festschrift for Curry. It develops the proofs-as-terms correspondence for implicational logic, extends it to the other propositional connectives, then to a term language for Heyting Arithmetic.
\end{itemize}
